\title{DeepSeekMath: Pushing the Limits of Mathematical Reasoning in Open Language Models}

\begin{abstract}
Mathematical reasoning remains a formidable obstacle for language models, largely due to the intricate and highly structured characteristics inherent in mathematics. This work presents DeepSeekMath~7B, which builds upon DeepSeek-Coder-Base-v1.5 7B through continued pre-training using 120B math-centric tokens collected from Common Crawl, in conjunction with textual and code data. DeepSeekMath~7B secures a notable 51.7\% on the challenging, competition-level MATH benchmark—accomplished without resorting to external toolkits or ensemble strategies—bringing its results close to the likes of Gemini-Ultra and GPT-4. Additionally, employing self-consistency across 64 sampled responses from DeepSeekMath~7B boosts performance to 60.9\% on MATH. The model’s aptitude in mathematical reasoning is primarily the result of two innovations: firstly, a carefully designed data selection pipeline is leveraged to maximize the utility of vast, publicly available web resources; secondly, Group Relative Policy Optimization (GRPO), a modified approach to Proximal Policy Optimization (PPO), is introduced to strengthen mathematical reasoning proficiency and simultaneously improve PPO's memory efficiency.
\end{abstract}

\section{Introduction}

The emergence of large language models (LLMs) has transformed strategies for tackling mathematical reasoning tasks in artificial intelligence, catalyzing substantial progress on both quantitative reasoning benchmarks \citep{MATH} and geometry reasoning benchmarks \citep{trinh2024solving}. Beyond benchmark improvements, these models have also served as powerful tools to aid humans with intricate mathematical problem-solving \citep{tao}. Nevertheless, state-of-the-art systems like GPT-4 \citep{gpt4} and Gemini-Ultra \citep{gemini} remain closed-source, and their public, open-source alternatives continue to underperform by a significant margin.

This paper presents DeepSeekMath, a specialized language model that demonstrates superior mathematical proficiency compared to open-source alternatives and nearly matches GPT-4’s performance on standard academic evaluations. The foundation for this model is the DeepSeekMath Corpus, an extensive, high-quality pre-training dataset consisting of 120B mathematical tokens. The construction of this dataset begins with extraction from the Common Crawl (CC), utilizing a classifier based on fastText \citep{joulin2016fasttext}. Initially, the classifier is trained with positive samples sourced from OpenWebMath \citep{openwebmath} and a wide variety of web pages serving as negative examples. The trained classifier is subsequently used to retrieve more positive samples from the CC, which are subjected to additional human annotation for quality assurance. The classifier is then retrained with these augmented data to enhance accuracy. Evaluations reveal that this large-scale corpus possesses high quality, evidenced by the performance of our base model, DeepSeekMath-Base 7B, which obtains 64.2\% on GSM8K \citep{gsm8k} and 36.2\% on the MATH competition dataset \citep{MATH}, surpassing Minerva 540B \citep{minerva}. Furthermore, as the DeepSeekMath Corpus incorporates multiple languages, improvements are observed on Chinese mathematics benchmarks \citep{wei2023cmath,agieval}. We consider our contributions to mathematical data curation to be a foundation for future research, highlighting substantial opportunities for further advancement.

DeepSeekMath-Base is initialized using DeepSeek-Coder-Base-v1.5 7B \citep{deepseek-coder}, as we determine that leveraging a code-focused pretraining model yields superior outcomes compared to beginning with a general-purpose LLM. Additionally, our findings suggest that math-specific pretraining not only enhances the model’s performance on mathematical tasks but also improves its results on broader benchmarks like MMLU \citep{mmlu} and BBH \citep{bbh}, demonstrating augmented general reasoning skills.

Subsequent to pretraining, we subject DeepSeekMath-Base to mathematical instruction tuning that utilizes data annotated with chain-of-thought \citep{cot}, program-of-thought \citep{pot,pal}, and tool-integrated reasoning \citep{tora} methodologies. This tuning yields DeepSeekMath-Instruct 7B, a model that surpasses all peer 7B models and achieves performance on par with much larger, 70B-scale open-source instruction-tuned models.

In addition, we propose Group Relative Policy Optimization (GRPO), an RL algorithm derived from Proximal Policy Optimization (PPO) \citep{schulman2017proximal}. GRPO eliminates the need for a critic model, instead utilizing group-based score baselines, which substantially decreases resource demands during training. Training on a targeted subset of English instruction data, GRPO delivers marked gains over the robust DeepSeekMath-Instruct baseline, across both in-domain tasks (GSM8K: 82.9\% $\rightarrow$ 88.2\%, MATH: 46.8\% $\rightarrow$ 51.7\%) and out-of-domain mathematical challenges (e.g., CMATH: 84.6\% $\rightarrow$ 88.8\%) during the RL phase. 

Moreover, we introduce a unified analytical framework that encapsulates multiple approaches—including Rejection Sampling Fine-Tuning (RFT) \citep{yuan2023scaling}, Direct Preference Optimization (DPO) \citep{dpo}, PPO, and GRPO—by conceptualizing these methods as forms of direct or simplified reinforcement learning. Leveraging this framework, we conduct comprehensive studies—such as comparisons between online and offline training, evaluations of outcome versus process supervision, and assessments of single-turn versus iterative RL—to elucidate the core contributors within this paradigm. Finally, we articulate the mechanisms through which our RL methodology enhances the performance of instruction-tuned models and outline future pathways for achieving more robust reinforcement learning under this unified conceptual structure.

\subsection{Contributions}
We present advancements in scalable mathematical pre-training, as well as an in-depth investigation into reinforcement learning methodologies.

\textbf{Math Pre-Training at Scale}

\begin{itemize}[topsep=0pt]
    \item Our study demonstrates that the widely available Common Crawl dataset possesses significant mathematical value. Leveraging a carefully engineered data selection pipeline, we have assembled the DeepSeekMath Corpus—a rigorously curated collection consisting of 120B tokens from web sources specifically filtered for mathematical relevance. This dataset is nearly 7 times larger than the mathematical portion utilized in Minerva \citep{minerva}, and approximately 9 times the scale of the recently introduced OpenWebMath dataset \citep{openwebmath}.

    \item DeepSeekMath-Base 7B, our foundational pre-trained model, performs on par with Minerva 540B \citep{minerva}, thereby illustrating that parameter count alone does not dictate mathematical reasoning performance. High-quality, domain-specific pre-training on a more modestly sized model can yield impressive results.

    \item We detail outcomes from our mathematical pre-training studies, revealing that preliminary code training notably enhances the model's proficiency in solving math problems, regardless of tool assistance. This finding contributes to the ongoing debate regarding the benefits of code pre-training, suggesting it indeed bolsters reasoning capabilities, particularly for mathematics.

    \item While it is standard practice to utilize arXiv papers in various math-focused research pipelines, our experiments indicate this approach does not deliver measurable gains across the mathematical evaluation tasks explored herein.
\end{itemize}

\textbf{Exploration and Analysis of Reinforcement Learning}

\begin{itemize}[topsep=0pt]
    \item We present Group Relative Policy Optimization (GRPO), a novel reinforcement learning approach that eliminates the need for a critic model. By estimating baselines via group scores, GRPO substantially reduces computational overhead during training when compared to Proximal Policy Optimization (PPO), while maintaining efficacy.
    \item Our results reveal that integrating GRPO into the instruction-tuning phase substantially improves the performance of the DeepSeekMath-Instruct model using only instruction-tuning datasets. Moreover, we observe improved generalization to out-of-domain tasks throughout the reinforcement learning stage.
    \item We articulate a cohesive framework for interpreting a variety of algorithms, encompassing RFT, DPO, PPO, and GRPO. Additionally, we carry out in-depth experiments—such as comparisons between online and offline training, outcome versus process-based supervision, and analyses contrasting single-turn with iterative reinforcement learning—to elucidate the foundational components underlying this unified view.
    \item Drawing from this overarching framework, we investigate the underlying mechanisms contributing to the efficacy of reinforcement learning approaches, and we outline several promising avenues for further enhancing reinforcement learning applied to LLMs.
\end{itemize}

\subsection{Summary of Evaluations and Metrics}

\begin{itemize}[topsep=0pt]
    \item \textbf{English and Chinese Mathematical Reasoning}: Our evaluation comprises rigorous testing on a range of English and Chinese benchmark datasets, encompassing math problems spanning from primary school through collegiate curricula. English datasets include GSM8K \citep{gsm8k}, MATH \citep{MATH}, SAT \citep{llemma}, OCW Courses \citep{minerva}, and MMLU-STEM \citep{mmlu}. For Chinese, we utilize MGSM-zh \citep{mgsm}, CMATH \citep{wei2023cmath}, Gaokao-MathCloze \citep{agieval}, and Gaokao-MathQA \citep{agieval}. Assessment focuses both on models’ capacities to generate comprehensive, tool-free solutions and to employ Python for problem-solving.

    On English-language tasks, DeepSeekMath-Base matches the closed-source Minerva 540B \citep{minerva} in capability and decisively outperforms all existing open-source base models (such as Mistral 7B \citep{mistral} and Llemma-34B \citep{llemma}), irrespective of whether they include math-specific pre-training. Significantly, DeepSeekMath-Base achieves even higher accuracy on Chinese evaluations, an advantage likely stemming from our broader pre-training data collection that, unlike previous efforts \citep{minerva,llemma}, incorporates substantial high-quality non-English data. After incorporating instruction-based tuning and reinforcement learning, the resulting DeepSeekMath-Instruct and DeepSeekMath-RL models demonstrate notable advances—most notably, surpassing the 50\% accuracy threshold on the challenging, competition-grade MATH dataset, a milestone first for open-source models. 
\end{itemize}

\item \textbf{Formal Mathematics}: To assess DeepSeekMath-Base in formal mathematical contexts, we employ the informal-to-formal theorem proving benchmark introduced by \citep{dsp_proof} on the miniF2F dataset \citep{minif2f}, utilizing Isabelle \citep{isabelle} as the selected proof assistant. Our results indicate that DeepSeekMath-Base achieves robust few-shot performance in autoformalization tasks.

\item \textbf{Natural Language Understanding, Reasoning, and Code}: For a broad evaluation of the model’s general understanding, reasoning, and programming proficiency, we test DeepSeekMath-Base on several benchmarks: the Massive Multitask Language Understanding (MMLU) benchmark \citep{mmlu} encompassing 57 diverse multiple-choice tasks, BIG-Bench Hard (BBH) \citep{bbh} with 23 multi-step reasoning problems, along with HumanEval \citep{codex} and MBPP \citep{mbpp}, both standard assessments of code generation ability. Results suggest that math-centric pre-training enhances not only reasoning but also natural language understanding skills.

\section{Math Pre-Training}

\subsection{Data Collection and Decontamination} This section details the methodology behind constructing the DeepSeekMath Corpus using data from Common Crawl. As shown in Figure \ref{fig:data_collect}, we adopt an iterative pipeline designed for efficient assembly of a large-scale mathematical text corpus, starting with a small, curated seed dataset (e.g., high-quality mathematical resources). Notably, this pipeline can be readily adapted for data collection in other domains such as programming.

Initially, we select OpenWebMath \citep{openwebmath}—a curated repository of mathematical texts on the web—as our seed dataset. With this data, we train a fastText classifier \citep{joulin2016fasttext} to identify additional math-related web pages with characteristics similar to those in OpenWebMath. Specifically, we randomly sample 500,000 entries from the seed as positive instances and 500,000 web pages from Common Crawl as negative samples. Training is carried out using an open-source toolkit\footnote{\url{https://fasttext.cc}}, setting the embedding dimension to 256, a learning rate of 0.1, a maximum word n-gram length of 3, a minimum word occurrence threshold of 3, and 3 training epochs. To manage the enormous scale of Common Crawl, we implement deduplication and near-duplicate filtering by URL, reducing the dataset to approximately 40B HTML pages. Next, we apply the fastText model to this deduplicated corpus to surface web pages likely to contain mathematical content. To further enhance quality, the retrieved pages are ranked by fastText-predicted scores, retaining only the highest-scoring entries. We conduct pre-training trials using the leading 40B, 80B, 120B, and 160B tokens, ultimately choosing the top 40B tokens in the initial iteration.

Following the initial round of data gathering, a significant number of mathematical web pages are still missing from the corpus. This gap is largely attributed to the limited variety within the positive examples used to train the fastText model. To address this, we broaden the seed corpus by sourcing additional mathematical web content, aiming to enhance the fastText model's capacity. The process begins by partitioning the entire Common Crawl dataset into distinct domains, with each domain comprising web pages that share a common base URL. For every domain, we compute the fraction of web pages collected during the first iteration. Domains in which more than 10\% of the pages are captured are designated as math-related (such as \url{mathoverflow.net}).

We then proceed to manually label the specific URLs within these math-related domains that pertain to mathematical content (for instance, \url{mathoverflow.net/questions}). Any previously uncollected pages associated with these URLs are subsequently incorporated into the seed set. By doing so, we can include more positive samples for training, allowing the fastText model to achieve higher recall of mathematical data in later rounds. After four cycles of data collection, we obtain a corpus of 35.5 million mathematical web pages, comprising a total of 120 billion tokens. Notably, during the fourth iteration, we observe that almost 98\% of the data had already been retrieved in the prior round, prompting us to halt further data collection at this point.

In order to mitigate the risk of benchmark data leakage, we adopt the filtering strategy used by \cite{deepseek-coder}: we exclude web pages that contain questions or answers originating from prominent mathematical benchmarks, including English resources such as GSM8K \citep{gsm8k} and MATH \citep{MATH}, as well as Chinese datasets like CMATH \citep{wei2023cmath} and AGIEval \citep{agieval}. Our filtering mechanism removes any text snippet in which a contiguous sequence of 10 tokens matches any substring found in the benchmark datasets. For benchmark entries shorter than 10 tokens but at least 3 tokens long, we apply an exact match filter to excise contaminated pages from the math corpus.

\subsection{Validating the Quality of the DeepSeekMath~Corpus}

We conduct a series of pre-training experiments to benchmark the DeepSeekMath~Corpus against newly published mathematical corpora:

\begin{itemize}[topsep=0pt]
    \item \textbf{MathPile} \citep{mathpile}: An 8.9B-token, multi-source collection that draws from textbooks, Wikipedia, ProofWiki, CommonCrawl, StackExchange, and arXiv, with arXiv contributing more than 85\% of the content;
    \item \textbf{OpenWebMath} \citep{openwebmath}: A 13.6B-token dataset curated from CommonCrawl by selecting for mathematical web content;
    \item \textbf{Proof-Pile-2} \citep{llemma}: A composite mathematical dataset that integrates OpenWebMath, AlgebraicStack (which offers 10.3B tokens of mathematical code), and arXiv manuscripts (28.0B tokens). For Proof-Pile-2 experiments, we maintain the arXiv:Web:Code distribution ratio of 2:4:1, as described in \cite{llemma}.
\end{itemize}

\subsubsection{Training Setting}
\label{sec:quality-policy}
We conduct math-specific training on a general-purpose language model comprising 1.3B parameters, which follows the architecture established by the DeepSeek LLMs \citep{deepseek-llm}, referred to as DeepSeek-LLM 1.3B. For each mathematical dataset, a distinct model is trained for 150B tokens. All training procedures utilize the streamlined HAI-LLM framework \citep{haillm}. Consistent with the protocols used for DeepSeek LLMs, we adopt the AdamW optimizer \citep{adamW} with parameters $\beta_1=0.9$, $\beta_2=0.95$, and $\mathrm{weight\_decay}=0.1$. The learning rate follows a multi-phase schedule: starting with 2,000 warmup steps to reach its maximum, subsequently reduced to 31.6\% at 80\% of the training duration, and further dropped to 10.0\% at the 90\% point. The maximum learning rate is set to 5.3e-4. During training, the batch size is 4M tokens, and the context length is fixed at 4K.

\subsubsection{Evaluation Results}

\textbf{The DeepSeekMath~Corpus stands out for its exceptional quality, extensive multilingual coverage in mathematics, and unprecedented scale.}

\begin{itemize}[topsep=0pt]
    \item \textbf{High-quality}:
    The assessment involves testing on 8 downstream mathematical benchmarks, employing few-shot chain-of-thought prompting \cite{cot}. Results in Table \ref{tab:corpora_comparison} indicate that models fine-tuned with the DeepSeekMath~Corpus outperform their counterparts. As illustrated in Figure \ref{fig:corpora_comparisons}, models trained on the DeepSeekMath Corpus surpass those trained on Proof-Pile-2 at 50B tokens (which corresponds to one complete pass through Proof-Pile-2), implying that the overall quality of DeepSeekMath Corpus is superior.
    \item \textbf{Multilingual}:
    The DeepSeekMath~Corpus contains mathematical texts in various languages, with English and Chinese constituting the largest shares. According to Table \ref{tab:corpora_comparison}, training on the DeepSeekMath~Corpus improves mathematical reasoning in both English and Chinese. By contrast, other mathematical datasets—primarily in English—do not provide significant gains and can even negatively impact Chinese mathematical reasoning abilities.
    \item \textbf{Large-scale}:
    The DeepSeekMath~Corpus is markedly larger than other mathematical corpora. Figure \ref{fig:corpora_comparisons} demonstrates that when DeepSeek-LLM 1.3B is trained with DeepSeekMath~Corpus, it exhibits a sharper and more sustained performance trajectory. In comparison, existing baseline corpora are considerably smaller and have been cycled through multiple times during training, which leads to an early stagnation in model performance.
\end{itemize}

\subsection{Training and Evaluating DeepSeekMath-Base 7B}

This section details the development of DeepSeekMath-Base 7B, a foundational model characterized by advanced mathematical reasoning skills. The model’s parameters are initialized using DeepSeek-Coder-Base-v1.5 7B \citep{deepseek-coder} and undergo pretraining over 500B tokens. The composition of the training dataset is distributed as follows: 56\% from the DeepSeekMath Corpus, 4\% from AlgebraicStack, 10\% sourced from arXiv, 20\% comprises Github code, and the last 10\% includes natural language text collected from Common Crawl in both English and Chinese. The primary training configuration mirrors the setup in Section \ref{sec:quality-policy}, with two key modifications: the learning rate is capped at a maximum of 4.2e-4, and a batch size of 10M tokens is employed.

To thoroughly evaluate the mathematical proficiency of DeepSeekMath-Base 7B, we analyze its performance in generating independent mathematical solutions without the use of auxiliary tools, leveraging tools for problem-solving, and engaging in formal theorem proving tasks. Additionally, the model is assessed across broader competencies, including its capabilities in natural language comprehension, logical reasoning, and programming.

\paragraph{Step-by-Step Mathematical Reasoning and Problem Solving}

The effectiveness of DeepSeekMath-Base in resolving mathematical questions is measured through few-shot chain-of-thought prompting \citep{cot}, tested on eight benchmarks available in both English and Chinese. These benchmarks include quantitative reasoning datasets such as GSM8K \citep{gsm8k}, MATH \citep{MATH}, and CMATH \citep{wei2023cmath}, as well as multiple-choice assessments like MMLU-STEM \citep{mmlu} and Gaokao-MathQA \citep{agieval}. The benchmarks collectively span a wide range of mathematical topics, from elementary to undergraduate-level difficulty.

As detailed in Table \ref{tab:base_math_cot}, DeepSeekMath-Base 7B consistently achieves the highest scores among open-source base models on all eight assessed benchmarks, outperforming established models such as Mistral 7B \citep{mistral} and the more recently introduced Llemma 34B \citep{llemma}, which benefits from specialized training on Proof-Pile-2 \citep{llemma}. Particularly noteworthy is its performance on the competition-grade MATH dataset, where DeepSeekMath-Base attains a margin exceeding 10\% absolute over other open-source models, and even surpasses Minerva 540B \citep{minerva}—a closed-source foundation model that is 77 times larger, derived from PaLM \citep{palm}, and further optimized with mathematical data.

\paragraph{Mathematical Problem Solving with Tool Use}  
For the assessment of program-based mathematical reasoning, we conduct evaluations on GSM8K and MATH by employing few-shot program-of-thought prompting strategies \citep{pot,pal}. Within this setup, models are instructed to devise a Python solution to each posed problem, making use of libraries such as \textit{math} and \textit{sympy} to handle sophisticated calculations. The computed output from executing these scripts is then used as the proposed answer. As reflected in Table \ref{tab:base_math_pot_proof}, DeepSeekMath-Base 7B achieves superior results, surpassing the previously leading Llemma 34B.

\paragraph{Formal Mathematics}

Automated formal proof generation offers both increased reliability and efficiency in producing mathematical arguments, and has attracted growing research interest recently. We benchmark DeepSeekMath-Base 7B using the informal-to-formal proof task presented in \citep{dsp_proof}, wherein the model is provided an informal statement, its corresponding formalization, and an informal proof, then tasked to output a formal proof. Our experiments are carried out on miniF2F \citep{minif2f}, which serves as a standard for Olympiad-level formal mathematics. For each test case, a formal proof is generated in Isabelle using few-shot prompting. Consistent with the protocol of \cite{dsp_proof}, model-generated proof outlines are further completed using the standard automated prover Sledgehammer \citep{sledgehammer} to supply detailed steps. According to the results in Table \ref{tab:base_math_pot_proof}, DeepSeekMath-Base 7B shows notable effectiveness in automatic proof formalization.

\paragraph{Natural Language Understanding, Reasoning, and Code}

We evaluate model capabilities in natural language understanding on MMLU \citep{mmlu}, logical reasoning on BBH \citep{bbh}, and programming skills on HumanEval \citep{codex} as well as MBPP \citep{mbpp}. The results displayed in Table \ref{tab:understanding_reasoning_code} reveal that DeepSeekMath-Base 7B achieves marked performance gains over the previous iteration, DeepSeek-Coder-Base-v1.5 \citep{deepseek-coder}, on both MMLU and BBH, underscoring the contribution of targeted math training to the improvement of reasoning and understanding tasks. Moreover, by including code token data during continued training, DeepSeekMath-Base 7B sustains the coding performance found in DeepSeek-Coder-Base-v1.5 on both programming benchmarks. Across all three benchmarks related to reasoning and code, DeepSeekMath-Base 7B substantially outperforms the general-purpose model Mistral 7B \citep{mistral}.

\section{Supervised Fine-Tuning}

\subsection{SFT Data Curation}

We assemble an instruction-tuning dataset for mathematics that includes both English and Chinese questions, covering a spectrum of mathematical disciplines and a range of difficulty. Each problem is associated with a detailed solution in chain-of-thought (CoT) \citep{cot}, program-of-thought (PoT) \citep{pot,pal}, or tool-integrated reasoning format \citep{tora}. The final dataset contains 776K training instances.

\begin{itemize}[topsep=0pt]
    \item \textbf{English mathematical datasets}: We supply tool-integrated solution annotations for GSM8K and MATH, and incorporate a selected portion of MathInstruct \citep{MathInstruct} together with the Lila-OOD training set \citep{lila}, all featuring CoT or PoT-style problem solving. This collection encompasses multiple mathematical areas, such as algebra, probability, number theory, calculus, and geometry.
    \item \textbf{Chinese mathematical datasets}: We compile K-12 mathematics questions in Chinese across 76 fine-grained categories, like linear equations, and annotate these with both CoT and tool-integrated reasoning solutions.
\end{itemize}

\subsection{Training and Evaluating DeepSeekMath-Instruct 7B}

Here, we describe the mathematical instruction tuning process for DeepSeekMath-Instruct 7B, using DeepSeekMath-Base as its foundation. Training data instances are concatenated at random up to a context limit of 4K tokens. The model undergoes 500 training iterations, using a batch size of 256 and a fixed learning rate of 5e-5.

Model performance on mathematics is assessed—both with and without the aid of tools—on four quantitative reasoning evaluation sets in English and Chinese. We perform comparative evaluation against the leading contemporary models:

\begin{itemize}[topsep=0pt]
    \item \textbf{Closed-source models} include: (1) the GPT series, notably GPT-4 \citep{gpt4} and GPT-4 Code Interpreter~\footnote{\url{https://openai.com/blog/chatgpt-plugins\#\#code-interpreter}}, (2) Gemini Ultra and Pro \citep{gemini}, (3) Inflection-2 \citep{inflection-2}, (4) Grok-1~\footnote{\url{https://x.ai/model-card}}, as well as several models recently released by Chinese technology firms, such as (5) Baichuan-3~\footnote{\url{https://www.baichuan-ai.com}} and (6) the latest GLM-4~\footnote{\url{https://open.bigmodel.cn/dev/api\#glm-4}} from the GLM series \citep{glm}.
\end{itemize}

These models serve general-purpose functions and typically have undergone various alignment strategies.
    \item \textbf{Open-source models} encompass both general LLMs and those with specific mathematical enhancements: (1) DeepSeek-LLM-Chat 67B \citep{deepseek-llm}, (2) Qwen 72B \citep{qwen}, (3) SeaLLM-v2 7B \citep{seallm}, and (4) ChatGLM3 6B \citep{chatglm3}. In addition, several models are tailored for mathematical tasks: (5) InternLM2-Math 20B~\footnote{\url{https://github.com/InternLM/InternLM-Math}}, which extends InternLM2 through additional math training and subsequent instruction tuning, (6) Math-Shepherd-Mistral 7B, which leverages PPO training \citep{schulman2017proximal} on Mistral 7B \citep{mistral} using a process-supervised reward model, (7) the WizardMath series \citep{wizardmath}, which advances mathematical reasoning capabilities in Mistral 7B and Llama-2 70B \citep{llama2} by employing evolve-instruct (a variant of instruction tuning utilizing AI-generated instructions) in combination with PPO on problems mainly from GSM8K and MATH, (8) MetaMath 70B \citep{metamath}, derived by fine-tuning Llama-2 70B on enriched versions of GSM8K and MATH, (9) ToRA 34B \cite{tora}, created by fine-tuning CodeLlama 34B for integrated tool-based mathematical reasoning, and (10) MAmmoTH 70B \citep{MathInstruct}, which builds on Llama-2 70B through instruction tuning with MathInstruct data.
\end{itemize}

Referencing Table \ref{tab:sft_rl_math}, when evaluated in settings prohibiting tool usage, DeepSeekMath-Instruct 7B displays strong stepwise reasoning performance. On the competitive MATH benchmark, it achieves a lead of at least 9\% absolute over all open-source contenders as well as most proprietary models (such as Inflection-2 and Gemini Pro). This performance margin persists even in comparison to much larger models (like Qwen 72B) and those trained specifically for mathematical reasoning using reinforcement learning techniques (for instance, WizardMath-v1.1 7B). Although DeepSeekMath-Instruct 7B matches the performance of leading Chinese proprietary models GLM-4 and Baichuan-3 on MATH, it is still outperformed by models such as GPT-4 and Gemini Ultra.

When models are allowed to incorporate both natural language reasoning and external programmatic tools during evaluation, DeepSeekMath-Instruct 7B attains nearly 60\% accuracy on the MATH dataset, outperforming all available open-source alternatives. For other benchmarks, its results are on par with DeepSeek-LLM-Chat 67B, the previous state-of-the-art model which is tenfold larger in size.

\section{Reinforcement Learning}

\subsection{Group Relative Policy Optimization} Reinforcement learning (RL) has demonstrated efficacy in further boosting the mathematical reasoning performance of LLMs subsequent to the Supervised Fine-Tuning (SFT) phase \citep{wang2023math,wizardmath}. Here, we present Group Relative Policy Optimization (GRPO), an RL method designed for efficiency and effectiveness.

\subsubsection{From PPO to GRPO} Proximal Policy Optimization (PPO) \citep{schulman2017proximal} is a widely adopted actor-critic RL approach in LLM fine-tuning via reinforcement learning \citep{ouyang2022training}. PPO refines LLM policies by maximizing the following surrogate objective:

\begin{equation}
\footnotesize
    \mathcal{J}_{PPO}(\theta) = \mathbb{E}{[q \sim P(Q), o \sim \pi_{\theta_{old}}(O|q)]} \frac{1}{|o|} \sum_{t=1}^{|o|} \min \left[ \frac{\pi_\theta(o_{t} | q, o_{<t})}{\pi_{\theta_{old}}(o_{t} | q, o_{<t})} A_{t}, \text{clip} \left( \frac{\pi_\theta(o_{t} | q, o_{<t})}{\pi_{\theta_{old}}(o_{t} | q, o_{<t})}, 1 - \varepsilon, 1 +

In this formulation, $\pi_{\theta}$ and $\pi_{\theta_{old}}$ denote the present and previous versions of the policy, respectively. The variables $q$ and $o$ correspond to questions sampled from the dataset and the outputs generated by the former policy $\pi_{\theta_{old}}$. The parameter $\varepsilon$ is a clipping hyper-parameter adopted in PPO to promote stability during training. The advantage term $A_t$ is derived using Generalized Advantage Estimation (GAE) \citep{gae}, relying on the sequence of rewards $\{r_{\ge t}\}$ along with a parameterized value function $V_{\psi}$. Consequently, PPO requires the simultaneous training of a value function together with the policy model. To guard against excessive reward model optimization, it is customary to include a KL penalty between the policy under training and a fixed reference model in the token-level reward at each generation step \citep{ouyang2022training}:

\begin{equation}
    r_{t} = r_\varphi(q, o_{\le t}) - \beta \log\frac{\pi_{\theta}(o_{t}|q, o_{<t})}{\pi_{ref}(o_{t}|q, o_{<t})},
\label{eq:PPO-reward}
\end{equation}

Here, $r_\varphi$ represents the reward function, $\pi_{ref}$ stands for the reference policy (commonly the initial SFT checkpoint), and $\beta$ is a scaling factor for the KL penalty.

Since the value function used by PPO often requires a model nearly as large as the policy model itself, this leads to significant demands on memory and computation. Moreover, the value function acts as a baseline in variance reduction when estimating the advantage during RL training. In large language model scenarios, however, the reward model generally assesses only the reward of the final token, which complicates learning a token-level value function with high fidelity. To mitigate these issues, and as illustrated in Figure \ref{fig:grpo}, we introduce Group Relative Policy Optimization (GRPO). GRPO eliminates the need to fit an explicit value function (as is done in PPO), instead employing the mean reward across a collection of outputs generated in response to the same question as the baseline. Specifically, given a question $q$, GRPO samples a set of responses $\{o_1, o_2, \cdots, o_G\}$ from $\pi_{\theta_{old}}$ and trains the policy model to optimize the objective below:

\begin{equation}
\footnotesize
\begin{split}
    \mathcal{J}_{GRPO}(\theta) &= \mathbb{E}{[q \sim P(Q), \{o_i\}_{i=1}^G \sim \pi_{\theta_{old}}(O|q)]}  \\
    & \frac{1}{G}\sum_{i=1}^G\frac{1}{|o_i|} \sum_{t=1}^{|o_i|} \left\{ \min \left[ \frac{\pi_\theta(o_{i,t} | q, o_{i,<t})}{\pi_{\theta_{old}}(o_{i,t} | q, o_{i,<t})} \hat{A}_{i,t}, \text{clip} \left( \frac{\pi_\theta(o_{i,t} | q, o_{i,<t})}{\pi_{\theta_{old}}(o_{i,t} | q, o_{i,<t})}, 1 - \varepsilon, 1 + \varepsilon \right)  \hat{A}_{i,t} \right] - \beta \mathbb{D}_{KL}\left[\pi_{\theta} || \pi_{ref}\right]\right\} ,
\end{split}
\label{eq:GRPO-obj}
\end{equation}

In this objective, $\varepsilon$ and $\beta$ serve as hyper-parameters. The term $\hat{A}_{i,t}$ refers to the group-relative advantage, computed solely based on intra-group reward differences—a process detailed in subsequent subsections. GRPO’s utilization of group-level relative advantage estimation aligns with the comparative methodology prevalent in reward modeling, which typically relies on pairwise output comparisons for individual questions. Additionally, in contrast to conventional methods where the KL penalty is integrated into the reward, GRPO imposes regularization by directly introducing the KL divergence between the current and reference policy into the loss function. This approach sidesteps any added complexity to the computation of $\hat{A}_{i,t}$. Unlike the per-step KL penalty in (\ref{eq:PPO-reward}), GRPO employs the following unbiased estimator for KL divergence as proposed

\begin{algorithm}[t]
  \small
  \caption{Iterative Group Relative Policy Optimization}
  \textbf{Input} initial policy model $\pi_{\theta_{\text{init}}}$; reward models $r_\varphi$; task prompts $\mathcal{D}$; 
  hyperparameters $\varepsilon$, $\beta$, $\mu$
  \begin{algorithmic}[1]
    \State Initialize the policy model $\pi_\theta$ as $\pi_{\theta_{\text{init}}}$
    \For{iteration = 1, \dots, I}
       \State Set the reference model $\pi_{ref}$ to the current $\pi_{\theta}$
      \For{step = 1, \dots, M}
      \State Randomly select a batch $\mathcal{D}_b$ from $\mathcal{D}$
      \State Assign the previous policy model $\pi_{\theta_{old}} \leftarrow \pi_{\theta}$ 
      \State For every question $q \in \mathcal{D}_b$, generate $G$ outputs $\{o_i\}_{i=1}^G \sim \pi_{\theta_{old}} (\cdot \mid q) $
      \State For each $o_i$, obtain rewards $\{r_i\}_{i=1}^{G}$ via the reward model $r_\varphi$
      \State For the $t$-th token in $o_i$, compute $\hat{A}_{i,t}$ using the group relative advantage estimator.
      \For{GRPO iteration = 1, \dots, $\mu$}
        \State Update $\pi_{\theta}$ by maximizing the GRPO objective function in Equation \ref{eq:GC-GRPO}
      \EndFor
    \EndFor \State Continuously train $r_\varphi$ by employing a replay buffer. 
    \EndFor 
  \end{algorithmic}
  \textbf{Output} $\pi_\theta$
  \label{alg:iter-grpo}
\end{algorithm}

\subsubsection{Outcome Supervision RL with GRPO} Formally, for each prompt $q$, a collection of $G$ outputs $\{o_1, o_2, \cdots, o_G\}$ is sampled from the previous policy $\pi_{\theta_{old}}$. These candidate outputs are then evaluated by a reward model, which produces the set of scores $\mathbf{r} = \{r_1, r_2, \cdots, r_G\}$. Each reward is standardized by subtracting the mean reward in the group and dividing by its standard deviation. In outcome supervision, the normalized reward is assigned at the final token of $o_i$, and the normalized value is also set as the advantage $\hat{A}_{i, t}$ for every token in $o_i$, specifically, $\hat{A}_{i, t} = \widetilde{r}_i = \frac{r_i- {\rm mean}(\mathbf{r})}{{\rm std}(\mathbf{r})}$. The policy is subsequently optimized by maximizing the objective described in equation (\ref{eq:GRPO-obj}).

\subsubsection{Process Supervision RL with GRPO} Since outcome supervision restricts feedback to only the output's end, it can be inadequate and inefficient for guiding policies, especially for intricate mathematical reasoning. Inspired by \cite{wang2023math}, we incorporate process supervision, wherein feedback is delivered after each reasoning step. Concretely, for a prompt $q$ and $G$ generated outputs $\{o_1, o_2, \cdots, o_G\}$, a process reward model scores each reasoning step. The resulting reward structure is $\mathbf{R} = \{ \{r_1^{index(1)},\cdots,r_1^{index(K_1)}\}, \cdots,  \{r_G^{index(1)},\cdots,r_G^{index(K_G)}\} \}$, where $index(j)$ marks the concluding token of step $j$ and $K_i$ counts steps in $o_i$. These stepwise rewards are normalized as $\widetilde{r}_i^{index(j)} = \frac{r_i^{index(j)} - {\rm mean(\mathbf{R})}}{{\rm std(\mathbf{R})}}$. Process supervision then computes the advantage for token $t$ by aggregating normalized rewards for all future steps, i.e., $\hat{A}_{i, t} = \sum_{index(j) \ge t} \widetilde{r}_i^{index(j)}$. The policy is updated by maximizing the objective from equation (\ref{eq:GRPO-

\subsubsection{Iterative RL with GRPO} During the reinforcement learning phase, the previously trained reward model may eventually fail to provide adequate supervision for the increasingly advanced policy model. To address this limitation, we implement an iterative RL framework leveraging GRPO. As detailed in Algorithm \ref{alg:iter-grpo}, the iterative approach involves collecting new datasets for the reward model through samples generated by the current policy model, and continuously refining the existing reward model using a replay strategy that retains 10\% of earlier data. Subsequently, we replace the reference model with the most recent policy model, and update the policy model in tandem with the newly improved reward model.

\subsection{Training and Evaluating DeepSeekMath-RL}

We apply reinforcement learning to DeepSeekMath-Instruct 7B, utilizing a set of approximately 144K chain-of-thought-format questions from GSM8K and MATH contained in the SFT data as training material. To examine the influence of RL on benchmarks where training examples are unavailable during the RL stage, we omit other SFT datasets. The construction of the reward model’s training set is aligned with the protocol described by \citep{wang2023math}. The initial reward model is trained atop DeepSeekMath-Base 7B with a learning rate of 2e-5. In GRPO, the policy model adopts a learning rate of 1e-6, and the KL divergence regularization coefficient is fixed at 0.04. For each query, $64$ candidate responses are sampled. The sequence length cap is 1024, and batch size during training is also 1024. Only a single policy model update is conducted for every exploration iteration. Evaluation of DeepSeekMath-RL 7B proceeds on the same benchmarks as DeepSeekMath-Instruct 7B. For this evaluation, chain-of-thought datasets like GSM8K and MATH are classified as in-domain, while all remaining benchmarks are treated as out-of-domain.

Table \ref{tab:sft_rl_math} presents a comparative analysis of open-source and closed-source models employing both chain-of-thought and tool-augmented reasoning on benchmark datasets in English and Chinese. Our key observations are as follows: 1) Using chain-of-thought reasoning, DeepSeekMath-RL 7B achieves accuracy rates of 88.2\% on GSM8K and 51.7\% on MATH, outperforming all open-source competitors within the 7B to 70B parameter range and surpassing most closed-source alternatives as well. 2) Notably, DeepSeekMath-RL 7B’s training data is exclusively limited to chain-of-thought-formatted instruction data from GSM8K and MATH, with training initiated from DeepSeekMath-Instruct 7B. Despite this restricted data regime, DeepSeekMath-RL 7B consistently exceeds the performance of DeepSeekMath-Instruct 7B across all evaluation categories, providing strong evidence for the efficacy of reinforcement learning.

\section{Discussion}
This section is dedicated to presenting insights derived from both pre-training and RL-based experimentation.

\subsection{Lessons Learnt in Pre-Training}

We begin by summarizing our experiences in the pre-training phase. Unless specified otherwise, all training conditions correspond to those detailed in Section \ref{sec:quality-policy}. Additionally, for references to the DeepSeekMath Corpus in this context, we rely on an 89-billion-token dataset compiled during the second iteration of data gathering.

\subsubsection{Code Training Benefits Mathematical Reasoning}

A prevailing yet largely anecdotal belief in the field is that pre-training on code enhances reasoning capabilities. Our investigation aims to shed light on this claim, particularly within the scope of mathematical reasoning: code training improves the model’s performance on mathematical reasoning tasks, regardless of the involvement of tool use.

To systematically examine the impact of code pre-training on mathematical reasoning, we evaluated two approaches: a two-stage training process and a one-stage training process:

\textbf{Two-Stage Training}

\begin{itemize}[topsep=0pt]
    \item \textbf{Code Training for 400B Tokens $\rightarrow$ Math Training for 150B Tokens}: DeepSeek-LLM 1.3B undergoes pretraining on 400B tokens from code corpora, after which it is further trained on 150B tokens consisting of mathematical content;
    \item \textbf{General Training for 400B Tokens $\rightarrow$ Math Training for 150B Tokens}: To serve as a baseline, we substitute the code data in the initial phase with general-purpose tokens (drawn from an expansive general text corpus curated by DeepSeek-AI) and repeat the two-stage training, aiming to systematically assess the comparative impact of code versus general text pretraining on downstream mathematical reasoning.
\end{itemize}

\textbf{One-Stage Training}

\begin{itemize}[topsep=0pt]
    \item \textbf{Math Training for 150B Tokens}: DeepSeek-LLM 1.3B is directly exposed to 150B math-specific tokens throughout its training process;
    \item \textbf{Training on a mixture of 400B Code Tokens and 150B Math Tokens}: Observing that math finetuning following prior code training can impair code proficiency, we further explore the effectiveness of interleaving code and math tokens within a single-stage training regime, evaluating whether this mitigates catastrophic forgetting while simultaneously enhancing mathematical reasoning skills.
\end{itemize}

\paragraph{Results}
Tables \ref{tab:code-math} and \ref{tab:code-math-general-eval} present evaluations on downstream tasks under the abovementioned experimental configurations.

Pretraining on code significantly promotes mathematical reasoning with programmatic assistance, a benefit observable in both two-stage and one-stage paradigms. Notably, Table \ref{tab:code-math} indicates that sequential code pretraining substantially boosts the model's competence on GSM8K and MATH benchmarks when leveraging Python for solutions. Introducing a subsequent math training stage leads to additional improvements. Under the one-stage strategy, jointly utilizing code and math tokens curtails the detrimental effects of catastrophic forgetting inherent in sequential learning, and further delivers enhancements in both code-oriented evaluation (Table \ref{tab:code-math-general-eval}) and math tasks with program support (Table \ref{tab:code-math}).

Moreover, code pretraining confers advantages even when mathematical reasoning occurs without external tool use. With two-stage training, the initial focus on code leads to measurable gains, and primes the model for more effective subsequent math-centric learning, resulting in peak task performance. In contrast, when code and math data are merged from the outset in a single-stage training scenario, the model’s performance on mathematical reasoning without tools diminishes. This outcome suggests that the limited capacity of DeepSeek-LLM 1.3B may restrict its ability to concurrently assimilate large volumes of both code and mathematical examples.

\subsubsection{ArXiv Papers Show Limited Effectiveness for Mathematical Reasoning Enhancement}

Although arXiv papers are widely integrated into mathematical pre-training datasets \citep{minerva,gpt-f,llemma,mathpile}, their direct contribution to enhancing mathematical reasoning abilities has not been thoroughly scrutinized. Interestingly, our empirical findings indicate that including arXiv papers does not lead to meaningful improvements in mathematical reasoning performance. To investigate this, we conducted experiments using models of varying scales, specifically DeepSeek-LLM 1.3B and DeepSeek-Coder-Base-v1.5 7B \citep{deepseek-coder}, training them on different arXiv-based corpora that had undergone distinct data processing procedures:

\begin{itemize}[topsep=0pt]
    \item \textbf{MathPile} \citep{mathpile}: a dataset containing 8.9B tokens, heavily cleaned and filtered via heuristic rules, with more than 85\% drawn from scientific arXiv papers;
    \item \textbf{ArXiv-RedPajama} \citep{redpajama}: a comprehensive collection of arXiv LaTeX documents from which preambles, comments, macros, and bibliographies have been excised, comprising a total of 28.0B tokens.
\end{itemize}

For these evaluations, DeepSeek-LLM 1.3B was pre-trained for 150B tokens and DeepSeek-Coder-Base-v1.5 7B for 40B tokens, each using one of the arXiv-focused corpora. Across the board, both models demonstrated little to no improvement, and at times even worse results, on a spectrum of mathematical reasoning benchmarks of varied difficulty featured in our study. The tasks evaluated encompassed quantitative reasoning datasets like GSM8K and MATH (refer to Table \ref{tab:arxiv-cot}), multiple-choice assessments such as MMLU-STEM (Table \ref{tab:arxiv-cot}), as well as formal mathematics problems like those in miniF2F (Table \ref{tab:arxiv_atp}).

Nevertheless, these findings must be considered within their methodological confines. Important open questions remain, such as:

\begin{itemize}[topsep=0pt]
    \item The specific influence of arXiv-derived data on math-related tasks not explored here, for example, informalization of theorems, where formal mathematical statements or proofs are translated into informal language;
    \item The outcomes when arXiv materials are supplemented with data from other domains;
    \item The possibility that larger-scale models could exhibit measurable benefits from arXiv pre-training.
\end{itemize}

Addressing these questions will require further research, which we propose for future work.

\subsection{Insights from Reinforcement Learning}
\subsubsection{Toward a Unified Training Paradigm} This section develops a unified framework for interpreting a variety of training approaches—including SFT, RFT, DPO, PPO, and GRPO—and empirically investigates which elements of this unified model are influential. Typically, the gradient of any such training strategy with respect to its parameter $\theta$ can be generally formulated as follows:

\begin{equation}
    \nabla_{\theta}\mathcal{J}_{\textcolor{red}{\mathcal{A}}}(\theta) = \mathbb{E}[\underbrace{(q,o) \sim \textcolor{red}{\mathcal{D}}}_{Data \ Source}]\left( \frac{1}{|o|} \sum_{t=1}^{|o|}  \underbrace{GC_{{\mathcal{A}}}(q, o, t, \textcolor{red}{\pi_{{rf}}})}_{Gradient \ Coefficient}  \nabla_{\theta}\log \pi_{\theta}(o_t | q, o_{<t})\right).
\label{eq:objective}
\end{equation}

This framework is structured around three principal elements: (1) the \textit{Data Source $\mathcal{D}$}, responsible for providing the dataset used for model training; (2) the \textit{Reward Function $\pi_{{rf}}$}, which serves as the origin of the feedback signal guiding learning; and (3) the \textit{Algorithm $\mathcal{A}$}, which integrates training samples and reward information to compute the gradient coefficient $GC$, thereby influencing the learning updates, either penalizing or reinforcing particular behaviors. The following summarizes typical approaches under this generalized scheme:

\begin{itemize}
    \item \textbf{Supervised Fine-tuning (SFT)}: SFT involves adjusting a pretrained model by further training it on curated datasets that have been manually selected.
    \item \textbf{Rejection Sampling Fine-tuning (RFT)}: RFT enhances the SFT-tuned model by additional fine-tuning using responses generated by the SFT model, filtered for answer correctness in response to SFT prompts.
    \item \textbf{Direct Preference Optimization (DPO)}: DPO builds upon the SFT model by applying fine-tuning with augmented candidate responses sampled from the SFT model, employing a pairwise loss based on preferences.
    \item \textbf{Online Rejection Sampling Fine-tuning (Online RFT)}: In contrast to RFT, Online RFT updates the policy using outputs sampled on-the-fly from the current policy model, starting from the SFT initialization, and continues fine-tuning using this dynamic data.
    \item \textbf{PPO/GRPO}: In these methods, the SFT model serves as the initial policy, which is then iteratively optimized using outputs sampled from the current (real-time) policy, guided by reinforcement learning principles.
\end{itemize}

Table \ref{tab:data_policy} provides an overview of the key elements constituting each method. For an expanded discussion and a step-by-step derivation, readers are referred to Appendix \ref{app:analysis-rl}.

\paragraph{Insight Regarding Data Source} The categorization of data sources distinguishes between online and offline sampling approaches. Online sampling involves collecting training data directly from the policy model as it is being trained in real time, whereas offline sampling utilizes data obtained from the outputs of a static, initial SFT model. Notably, RFT and DPO are characterized by their use of offline sampling, while Online RFT and GRPO rely on online sampling methodologies.

Figure \ref{fig:rl-analysis} illustrates a clear trend: Online RFT delivers marked improvements over RFT on both evaluated benchmarks. In the early training phase, Online RFT achieves performance levels similar to those of RFT. However, as training progresses, Online RFT secures a pronounced performance lead, underscoring the merits of online sampling. This outcome aligns with expectations: Initially, both the policy actor and the SFT model produce closely aligned outputs, and thus the sampled data reflects minimal variation. Over time, as the actor diverges, the data sourced from it exhibits greater distinctions, thereby amplifying the benefits conferred by dynamic, real-time data collection.

\paragraph{Insight Regarding Gradient Coefficient} The algorithms convert reward signals into gradient coefficients to facilitate model parameter updates. In our experimental setup, we delineate the reward functions as either `Rule' or `Model'. A `Rule'-based approach entails evaluating the correctness of responses against predefined criteria, while the `Model'-based approach employs a learned reward model that assigns scores to responses, where the reward model itself is trained using rule-based assessments. Equations \ref{eq:GC-RFT} and \ref{eq:GC-GRPO} clarify a fundamental distinction: GRPO introduces a mechanism for adjusting the gradient coefficient in proportion to the reward values issued by the reward model, thereby supporting fine-grained reinforcement or penalization dependent on response quality. In contrast, Online RFT does not differentiate among responses in this way; it lacks any penalty for incorrect answers and applies uniform positive reinforcement to all responses identified as correct.

As depicted in Figure \ref{fig:rl-analysis}, GRPO outperforms online RFT, underscoring the advantages of modifying the coefficients for positive and negative gradients. Moreover, the results for GRPO+PS exceed those for GRPO+OS, illustrating that leveraging more granular, step-dependent gradient coefficients provides notable benefits. We further investigate the effect of iterative RL by conducting two iterations in our experiments. Figure \ref{fig:iter-rl} shows that iterative RL considerably boosts performance, with the most substantial gains observed during the initial iteration.

\subsubsection{Why RL Works?} This work applies reinforcement learning to a select portion of the instruction tuning dataset, resulting in marked improvements over the instruction-tuned model. To delve deeper into the effectiveness of reinforcement learning, we assess the Pass@K and Maj@K accuracies for both the Instruct and RL models on two evaluation benchmarks. Figure \ref{fig:maj-pass} illustrates that RL raises Maj@K performance, yet leaves Pass@K largely unchanged. These results suggest that reinforcement learning principally strengthens the resilience of the model’s output distribution, i.e., \textbf{the gains appear to arise from improving the chance that a correct answer is among the TopK outputs, rather than an upgrade in core competencies.} Likewise, \citep{wang2023making} noted a \textbf{misalignment problem} in reasoning tasks for SFT models, demonstrating that such models’ reasoning abilities can be enhanced by introducing preference alignment methods \citep{yuan2023rrhf,song2023preference,wang2023making}.

\subsubsection{How to Achieve More Effective RL?}
\label{sec:effec_rl}
Our results confirm that RL is highly effective for mathematical reasoning tasks. In addition, we present a comprehensive framework to conceptualize and unify various prominent training strategies, treating each as a particular instance of direct or simplified RL methods. Equation \ref{eq:objective} encapsulates three central elements: Data Source, Algorithm, and Reward Function. We offer perspectives on possible research trajectories concerning each aspect.

\paragraph{Data Source} The data source underpins all training procedures. Within RL, this refers specifically to unlabeled queries, paired with responses sampled from the policy model. In this study, our RL framework utilizes only prompts originating from the instruction tuning dataset, employing straightforward nucleus sampling for output generation. We posit that this might explain why our RL approach yields improvements solely in Maj@K metrics. Future work will investigate applying our RL scheme to out-of-distribution prompts, possibly combined with \textbf{advanced sampling (decoding) strategies} such as those based on tree-search algorithms \citep{yao2023tree}. Furthermore, \textbf{efficient inference techniques} \citep{xia-etal-2023-speculative,leviathan2023fast,kwon2023efficient,xia2024unlocking}, which impact how effectively the policy model explores, are anticipated to be critically influential.

\paragraph{Algorithms} Algorithms utilize both input data and reward signals to compute gradient coefficients, thereby adjusting the model parameters. According to Equation \ref{eq:objective}, current methodologies predominantly \textbf{TRUST} the reward function’s feedback to modify the conditional probabilities of particular tokens—either increasing or decreasing them accordingly. Nonetheless, it cannot be guaranteed that the reward signal is consistently accurate, especially in the context of highly intricate tasks. As an illustration, even the meticulously curated PRM800K dataset \citep{lightman2023let}, which is annotated by skilled annotators, reportedly includes about 20\% incorrect annotations\footnote{\url{https://github.com/openai/prm800k/issues/12\#issuecomment-1728491852}}. Given these limitations, we aim to investigate reinforcement learning algorithms that remain robust in the presence of noisy reward signals. We posit that \textbf{WEAK-TO-STRONG} alignment approaches \citep{burns2023weak} have the potential to drive significant innovation in algorithmic design.

\paragraph{Reward Function} The reward function provides the principal learning signal throughout training. In the context of reinforcement learning, this function is commonly realized as a neural reward model. We identify three pivotal areas of development for these reward models: 1) \textbf{Enhancing the generalization capabilities of the reward model.} Reward models must generalize beyond the training distribution to competently assess out-of-distribution prompts and novel decoding outcomes. Failure to achieve this risks having reinforcement learning simply entrench the pre-existing distribution of large language models (LLMs) rather than advance their underlying capacities; 2) \textbf{Capturing and expressing the uncertainty of the reward model.} Representing model uncertainty may serve as a critical intermediary between weak reward models and weak-to-strong learning strategies; 3) \textbf{Developing high-fidelity process reward models efficiently} to yield detailed training feedback during stepwise reasoning processes \citep{lightman2023let,wang2023math}.

\section{Conclusion, Limitation, and Future Work}

We introduce DeepSeekMath, an open-source model that establishes new state-of-the-art results among open-source systems on the MATH benchmark at the competition level, achieving performance that rivals that of closed-source alternatives. Built on DeepSeek-Coder-v1.5 7B as its base, DeepSeekMath is continually trained over 500B tokens, with a notable emphasis on mathematical content, comprising 120B math-specific tokens gathered from Common Crawl. Through a comprehensive ablation analysis, we find that mathematical web pages constitute a promising resource for extracting quality mathematical data, while the utility of arXiv appears to be more limited than initially assumed. We further present Group Relative Policy Optimization (GRPO)—an adaptation of Proximal Policy Optimization (PPO)—which enhances mathematical reasoning skills while requiring less memory. Our results indicate that GRPO consistently boosts performance, even when DeepSeekMath-Instruct 7B already attains high benchmark scores. We also put forward a unified conceptual framework for interpreting multiple methodologies and delineate various prospective avenues for advancing reinforcement learning efficiency.

Despite DeepSeekMath's strong results on quantitative reasoning assessments, it underperforms compared to closed models in geometry and theorem-proving tasks. Specifically, during validation, the model struggled with questions involving triangles and ellipses, which may reflect limitations and biases introduced during data selection for pre-training and fine-tuning stages. Additionally, due to its model size, DeepSeekMath lags behind GPT-4 in terms of few-shot learning abilities: whereas GPT-4 demonstrates improvement with few-shot prompts, DeepSeekMath maintains similar scores between zero-shot and few-shot settings. Going forward, we plan to refine our engineered data curation strategy to assemble even higher-quality pre-training datasets. Moreover, we aim to pursue the potential strategies outlined in Section~\ref{sec:effec_rl} to further enhance reinforcement learning in large language models.